\begin{textsample}{POS Dim 1 – vem\_ed\_01 – Score 22.00 – vem\_ed\_01/t003.txt}  \label{ex:f1_pos_093}
Força de um grupo de pesquisa No primeiro semestre de 2020, sete Fatecs participam do PCI"Learning English as an additional language", com a universidade chinesa Tianjin Normal University. \textbf{São} as unidades de Franca, Americana, São José do Rio Preto, Itaquaquecetuba, Jahu, Piracicaba e Sebrae (em São Paulo).
A \textbf{partir} de 23 de março, as atividades \textbf{foram} pausadas, retornando no fim de abril.
A primeira colaboração ocorreu no segundo semestre de 2018; portanto, este já \textbf{é} o quarto semestre consecutivo de trabalhos com a instituição chinesa.
Neusa Haruka Sezaki Gritti, da Fatec Itaquaquecetuba, relata o \textbf{processo} das \textbf{interações}:"\textbf{Cada} equipe \textbf{é} composta de um professor de \textbf{cada} país (Brasil e China) e as \textbf{turmas} \textbf{são} organizadas em grupos de 6 a 8 alunos brasileiros e chineses, \textbf{que} se comunicam via WeChat ou GroupMe para explorar o trabalho."Sempre surgem temas levantados pelos estudantes e o contato online ocorre em videoconferências pelo Zoom, \textbf{ideal} para desenvolver a oralidade.
O trabalho pronto \textbf{segue} para a ferramenta Edmodo.
Segundo Mara Miranda, da Fatec Piracicaba, \textbf{é}"uma oportunidade para nossos estudantes interagirem, de \textbf{forma} direcionada pelos professores, com \textbf{pessoas} não falantes do Português".
Essas experiências contribuem para a profissionalização internacional dos alunos das Fatecs.
Paula Pudo entrou para a equipe no segundo semestre de 2019 e se considera"uma novata".
Ela leciona inglês nas Fatecs Itaquaquecetuba e Mogi das Cruzes e acredita \textbf{que} a experiência"ajuda e muito os alunos, devido ao \textbf{desenvolvimento} e aprendizagem do idioma-alvo, a língua inglesa".
Outras \textbf{contribuições}, no entender da professora, estão"no aprimoramento de \textbf{habilidades} \textbf{linguísticas} e comunicativas, além de \textbf{ampliar} os \textbf{horizontes} culturais."\textbf{Coordenação} proativa Gestores \textbf{que} vislumbram potencial dos \textbf{intercâmbios} virtuais mobilizam suas equipes para \textbf{criar} PCIs de sucesso. \textbf{É} o caso de Ricardo Sérgio Neiva Nóbrega, coordenador do curso de Comércio Exterior e professor de Gestão e Estratégia Internacional na Fatec Indaiatuba.
Ele lidera o projeto"Desenvolvendo internacionais", desenvolvido com Mona Pearl da DePaul University (EUA).
O PCI envolve grupos mistos de brasileiros (concluintes de Comércio Exterior) e norte-americanos.
Juntos, \textbf{identificam} empresas reais nos Estados Unidos e no Brasil e negociam uma aliança para se internacionalizarem em um \textbf{terceiro} país.
A colaboração \textbf{começou} no segundo semestre de 2019; em \textbf{cada} edição estiveram envolvidos cerca de 50 alunos (25 de \textbf{cada} país) e cinco professores, \textbf{sendo} um da DePaul e quatro da Fatec Indaiatuba: além de Nóbrega, participam Renata Pierri Lucietto Rossetto (Marketing Internacional), Carlos Antonio de Lima Penhalber (Negócios Internacionais) e Talita Annunciato Rodrigues (Inglês).
Segundo Nóbrega, a iniciativa ajuda os \textbf{futuros} tecnólogos a desenvolver"\textbf{competências} interculturais e interpessoais, \textbf{capacidade} de trabalhar em equipes mistas, resolver problemas e usar tecnologias".

% matched lemmas: ampliar, cada, capacidade, começar, competência, contribuição, coordenação, criar, desenvolvimento, forma, futuro, habilidade, horizonte, ideal, identificar, interação, intercâmbio, linguístico, partir, pessoa, processo, que, seguir, ser, terceiro, turma
\end{textsample}
